% =====================================================================
% Cine-JEPA — Results
% NeurIPS 2026 submission
% Standalone .tex fragment: include via % =====================================================================
% Cine-JEPA — Results
% NeurIPS 2026 submission
% Standalone .tex fragment: include via % =====================================================================
% Cine-JEPA — Results
% NeurIPS 2026 submission
% Standalone .tex fragment: include via % =====================================================================
% Cine-JEPA — Results
% NeurIPS 2026 submission
% Standalone .tex fragment: include via \input{results} from the main file.
% Citation keys are placeholders; resolve them in the bibliography.
% =====================================================================

\section{Results}
\label{sec:results}

We report decoding performance under the protocol of
\S\ref{sec:method:eval}: every method exposes a deterministic
clip-to-feature map and is fed into the same closed-form linear
probe on the same R6 test split. Each headline number is built in
two stages. \emph{Within} each pretraining seed we compute a
recording-level bootstrap mean of the test metric: the $108$ R6
recordings are resampled with replacement $B = 2{,}000$ times,
the metric is recomputed on each resample, and we take the
bootstrap mean as the per-seed score. \emph{Across} pretraining
seeds we then report the mean and one standard deviation of the
per-seed bootstrap means over five encoder seeds (CBraMod: four,
because of compute constraints). A row is marked
``$\checkmark$'' when the per-seed bootstrap means beat their
chance value at $p < 0.05$ under a one-sample two-sided $t$-test
($r = 0$ for Pearson $r$; $0.5$ for AUC and balanced accuracy;
$0.05/0.25$ for top-$1$/top-$5$ over $20$ bins).

% =====================================================================
\subsection{Stimulus-locked decoding}
\label{sec:results:stim}

Table~\ref{tab:stim-decoding} reports Pearson $r$ between probe
predictions and four continuous movie-feature targets — frame
luminance, frame contrast, normalized temporal position in the
movie, and the narrative event score — on the held-out R6 test
split.

\input{tables/tab_stim_decoding.tex}

\paragraph{Cine-JEPA tops every target.}
Cine-JEPA reaches Pearson $r = 0.226$ on luminance, $0.223$ on
contrast, $0.226$ on temporal position, and $0.155$ on the
narrative event score — the highest score in every column. Its
margin over the strongest non-Cine-JEPA method ranges from $9\%$
relative on luminance (vs.\ supervised Deep4Net at $0.207$) to
$26\%$ on contrast (vs.\ CorrCA-stats at $0.177$) and $23\%$ on
temporal position (vs.\ Luna full-finetune at $0.184$); on
narrative it is the only method to clear $r = 0.1$.

\paragraph{Without a stimulus-locked spatial filter, linear features carry no signal.}
Seven simple summary statistics (mean, std, log-power in
$\delta\!/\!\theta\!/\!\alpha\!/\!\beta\!/\!\gamma$ bands)
computed on the raw $129$-electrode input are statistically
indistinguishable from chance on every target ($|r| \le 0.016$,
all $n.s.$). The same statistics computed on the five
CorrCA-projected channels jump to $r \approx 0.16$--$0.18$ on
luminance, contrast, and temporal position. CorrCA is therefore
doing most of the heavy lifting that any \emph{linear} method can
extract from the EEG. It is, however, far from sufficient:
Cine-JEPA improves on CorrCA-stats by $0.04$--$0.07$ on the three
low/mid-level targets and by an absolute $+0.115$ on the
narrative event score, so the encoder is contributing real
nonlinear structure on top of the spatial filter.

\paragraph{The standard naturalistic-EEG decoder lags every learned method.}
Linear TRFs — the field-standard tool for naturalistic EEG
decoding \citep{crosse2016mtrf} — top out at $r = 0.117$ on
luminance under the CorrCA-5 input and decay to chance on
narrative ($r = 0.038$, $n.s.$). The fully-supervised end-to-end
Deep4Net is much stronger ($r = 0.207, 0.155, 0.146, 0.144$
across the four targets) but still trails Cine-JEPA on every
target by $0.011$ to $0.068$.

\paragraph{Foundation-model probes plateau below Cine-JEPA.}
Among the EEG foundation models we evaluated, Luna is the
strongest: a frozen linear probe on Luna features reaches
$r = 0.182, 0.158, 0.171, 0.053$, second only to Cine-JEPA on
three of four targets. Full fine-tuning of Luna gives a small
additional gain on contrast and temporal position but does not
close the gap, and feeding Luna a CorrCA-5 input
\emph{degrades} its scores rather than improves them — Luna's
own input pipeline already extracts a richer projection than
five linear components, so collapsing the channel dimension up
front discards information. BIOT is intermediate, CBraMod is
within noise of chance under our protocol, and BENDR / LaBraM
did not produce usable scores under the protocol and are listed
as ``---''.

\paragraph{Narrative event decoding is the discriminating target.}
All methods take their largest hit on the narrative event score:
Cine-JEPA is the only method to clear $r = 0.1$, and the
runner-up (Deep4Net at $0.144$) is itself supervised end-to-end
on the same target. Because the narrative score is higher-level
and less stimulus-locked than per-frame luminance or contrast,
the gap on this column separates methods that have learned a
temporally structured representation from methods that have
mainly captured low-level visual statistics. We return to this
in the discussion.

% =====================================================================
\subsection{Subject-trait decoding}
\label{sec:results:traits}
\todo{Fill in once the age-regression and sex-classification
results are finalized. Headline framing: do the same frozen
representations that win at stimulus-locked decoding also carry
recoverable subject-level structure (age as continuous Pearson
$r$; sex as ROC-AUC and balanced accuracy)? Reference
Table~\ref{tab:trait-decoding}.}

% =====================================================================
\subsection{Movie identity}
\label{sec:results:movie-id}
\todo{Fill in once the $20$-bin top-$1$ / top-$5$ movie-identity
results are finalized. Headline framing: position-in-movie can be
decoded as a continuous target (Table~\ref{tab:stim-decoding});
binning it into $20$ classes tests whether the representation
encodes coarser temporal context as a categorical signal.
Reference Table~\ref{tab:movie-id}.}

% =====================================================================
\subsection{Ablations}
\label{sec:results:ablations}
\todo{Fill in. Recommended ablations to present here: (i) Raw
$C{=}129$ vs.\ CorrCA-5 input under the same Cine-JEPA pipeline;
(ii) per-recording vs.\ global Z-score; (iii) VICReg coefficients
$\lambda_\sigma = \lambda_c \in \{0, 0.25, 1.0\}$; (iv) predictor
bottleneck $D' \in \{16, 24, 32, 64\}$; (v) visual-processing
delay $\Delta$ sweep, validating the $300$\,ms choice from
\S\ref{sec:method:data}.}
 from the main file.
% Citation keys are placeholders; resolve them in the bibliography.
% =====================================================================

\section{Results}
\label{sec:results}

We report decoding performance under the protocol of
\S\ref{sec:method:eval}: every method exposes a deterministic
clip-to-feature map and is fed into the same closed-form linear
probe on the same R6 test split. Each headline number is built in
two stages. \emph{Within} each pretraining seed we compute a
recording-level bootstrap mean of the test metric: the $108$ R6
recordings are resampled with replacement $B = 2{,}000$ times,
the metric is recomputed on each resample, and we take the
bootstrap mean as the per-seed score. \emph{Across} pretraining
seeds we then report the mean and one standard deviation of the
per-seed bootstrap means over five encoder seeds (CBraMod: four,
because of compute constraints). A row is marked
``$\checkmark$'' when the per-seed bootstrap means beat their
chance value at $p < 0.05$ under a one-sample two-sided $t$-test
($r = 0$ for Pearson $r$; $0.5$ for AUC and balanced accuracy;
$0.05/0.25$ for top-$1$/top-$5$ over $20$ bins).

% =====================================================================
\subsection{Stimulus-locked decoding}
\label{sec:results:stim}

Table~\ref{tab:stim-decoding} reports Pearson $r$ between probe
predictions and four continuous movie-feature targets — frame
luminance, frame contrast, normalized temporal position in the
movie, and the narrative event score — on the held-out R6 test
split.

% Stimulus-locked decoding table — referenced from results.tex
% Requires: booktabs
\begin{table}[t]
\centering
\small
\caption{Stimulus-locked decoding on the held-out R6 test split.
All values are Pearson $r$, mean $\pm 1\sigma$ across $5$ encoder
seeds (CBraMod: $n{=}4$), after a recording-level bootstrap
($B{=}2{,}000$) within each seed. ``$\checkmark$'' marks rows
significant against $r{=}0$ at $p{<}0.05$ (one-sample two-sided
$t$-test on the per-seed bootstrap means); ``$n.s.$'' marks
non-significance. Best per column in \textbf{bold}.}
\label{tab:stim-decoding}
\begin{tabular}{lcccc}
\toprule
Method & luminance & contrast & position & narrative \\
\midrule
chance & $0.000$ & $0.000$ & $0.000$ & $0.000$ \\
\addlinespace[2pt]
\multicolumn{5}{l}{\textit{Trivial baselines (no encoder)}} \\
Raw $129$-ch stats & $-0.016 \pm 0.005\,n.s.$ & $-0.008 \pm 0.022\,n.s.$ & $+0.006 \pm 0.025\,n.s.$ & $-0.006 \pm 0.015\,n.s.$ \\
CorrCA stats       & $0.167 \pm 0.041\,\checkmark$ & $0.177 \pm 0.043\,\checkmark$ & $0.158 \pm 0.032\,\checkmark$ & $0.040 \pm 0.023\,\checkmark$ \\
\addlinespace[2pt]
\multicolumn{5}{l}{\textit{Linear / supervised baselines}} \\
mTRF ($129$ channels) & $0.052 \pm 0.033\,\checkmark$ & $0.044 \pm 0.040\,n.s.$ & $0.027 \pm 0.037\,n.s.$ & $0.032 \pm 0.014\,\checkmark$ \\
mTRF (CorrCA-$5$)     & $0.117 \pm 0.011\,\checkmark$ & $0.052 \pm 0.026\,\checkmark$ & $0.071 \pm 0.032\,\checkmark$ & $0.038 \pm 0.040\,n.s.$ \\
Supervised end-to-end (Deep4Net) & $0.207 \pm 0.013\,\checkmark$ & $0.155 \pm 0.017\,\checkmark$ & $0.146 \pm 0.030\,\checkmark$ & $0.144 \pm 0.025\,\checkmark$ \\
\addlinespace[2pt]
\multicolumn{5}{l}{\textit{EEG foundation models, linear probe (frozen)}} \\
BENDR  & --- & --- & --- & --- \\
BIOT   & $0.127 \pm 0.033\,\checkmark$ & $0.107 \pm 0.035\,\checkmark$ & $0.107 \pm 0.038\,\checkmark$ & $0.003 \pm 0.009\,n.s.$ \\
LaBraM & --- & --- & --- & --- \\
Luna   & $0.182 \pm 0.021\,\checkmark$ & $0.158 \pm 0.018\,\checkmark$ & $0.171 \pm 0.012\,\checkmark$ & $0.053 \pm 0.036\,\checkmark$ \\
\addlinespace[2pt]
\multicolumn{5}{l}{\textit{EEG foundation models, full fine-tune}} \\
BENDR  & --- & --- & --- & --- \\
BIOT   & $0.085 \pm 0.060\,\checkmark$ & $0.074 \pm 0.066\,n.s.$ & $0.109 \pm 0.030\,\checkmark$ & $0.013 \pm 0.028\,n.s.$ \\
LaBraM & --- & --- & --- & --- \\
Luna   & $0.192 \pm 0.120\,\checkmark$ & $0.169 \pm 0.038\,\checkmark$ & $0.184 \pm 0.031\,\checkmark$ & $0.095 \pm 0.100\,n.s.$ \\
Luna + CorrCA-$5$ & $0.159 \pm 0.035\,\checkmark$ & $0.134 \pm 0.059\,\checkmark$ & $0.133 \pm 0.062\,\checkmark$ & $0.058 \pm 0.040\,\checkmark$ \\
CBraMod ($n{=}4$) & $0.052 \pm 0.077\,n.s.$ & $0.034 \pm 0.042\,n.s.$ & $0.022 \pm 0.079\,n.s.$ & $0.040 \pm 0.067\,n.s.$ \\
\addlinespace[2pt]
\textbf{Cine-JEPA (ours)} & $\mathbf{0.226 \pm 0.018}\,\checkmark$ & $\mathbf{0.223 \pm 0.014}\,\checkmark$ & $\mathbf{0.226 \pm 0.016}\,\checkmark$ & $\mathbf{0.155 \pm 0.023}\,\checkmark$ \\
\bottomrule
\end{tabular}
\end{table}


\paragraph{Cine-JEPA tops every target.}
Cine-JEPA reaches Pearson $r = 0.226$ on luminance, $0.223$ on
contrast, $0.226$ on temporal position, and $0.155$ on the
narrative event score — the highest score in every column. Its
margin over the strongest non-Cine-JEPA method ranges from $9\%$
relative on luminance (vs.\ supervised Deep4Net at $0.207$) to
$26\%$ on contrast (vs.\ CorrCA-stats at $0.177$) and $23\%$ on
temporal position (vs.\ Luna full-finetune at $0.184$); on
narrative it is the only method to clear $r = 0.1$.

\paragraph{Without a stimulus-locked spatial filter, linear features carry no signal.}
Seven simple summary statistics (mean, std, log-power in
$\delta\!/\!\theta\!/\!\alpha\!/\!\beta\!/\!\gamma$ bands)
computed on the raw $129$-electrode input are statistically
indistinguishable from chance on every target ($|r| \le 0.016$,
all $n.s.$). The same statistics computed on the five
CorrCA-projected channels jump to $r \approx 0.16$--$0.18$ on
luminance, contrast, and temporal position. CorrCA is therefore
doing most of the heavy lifting that any \emph{linear} method can
extract from the EEG. It is, however, far from sufficient:
Cine-JEPA improves on CorrCA-stats by $0.04$--$0.07$ on the three
low/mid-level targets and by an absolute $+0.115$ on the
narrative event score, so the encoder is contributing real
nonlinear structure on top of the spatial filter.

\paragraph{The standard naturalistic-EEG decoder lags every learned method.}
Linear TRFs — the field-standard tool for naturalistic EEG
decoding \citep{crosse2016mtrf} — top out at $r = 0.117$ on
luminance under the CorrCA-5 input and decay to chance on
narrative ($r = 0.038$, $n.s.$). The fully-supervised end-to-end
Deep4Net is much stronger ($r = 0.207, 0.155, 0.146, 0.144$
across the four targets) but still trails Cine-JEPA on every
target by $0.011$ to $0.068$.

\paragraph{Foundation-model probes plateau below Cine-JEPA.}
Among the EEG foundation models we evaluated, Luna is the
strongest: a frozen linear probe on Luna features reaches
$r = 0.182, 0.158, 0.171, 0.053$, second only to Cine-JEPA on
three of four targets. Full fine-tuning of Luna gives a small
additional gain on contrast and temporal position but does not
close the gap, and feeding Luna a CorrCA-5 input
\emph{degrades} its scores rather than improves them — Luna's
own input pipeline already extracts a richer projection than
five linear components, so collapsing the channel dimension up
front discards information. BIOT is intermediate, CBraMod is
within noise of chance under our protocol, and BENDR / LaBraM
did not produce usable scores under the protocol and are listed
as ``---''.

\paragraph{Narrative event decoding is the discriminating target.}
All methods take their largest hit on the narrative event score:
Cine-JEPA is the only method to clear $r = 0.1$, and the
runner-up (Deep4Net at $0.144$) is itself supervised end-to-end
on the same target. Because the narrative score is higher-level
and less stimulus-locked than per-frame luminance or contrast,
the gap on this column separates methods that have learned a
temporally structured representation from methods that have
mainly captured low-level visual statistics. We return to this
in the discussion.

% =====================================================================
\subsection{Subject-trait decoding}
\label{sec:results:traits}
\todo{Fill in once the age-regression and sex-classification
results are finalized. Headline framing: do the same frozen
representations that win at stimulus-locked decoding also carry
recoverable subject-level structure (age as continuous Pearson
$r$; sex as ROC-AUC and balanced accuracy)? Reference
Table~\ref{tab:trait-decoding}.}

% =====================================================================
\subsection{Movie identity}
\label{sec:results:movie-id}
\todo{Fill in once the $20$-bin top-$1$ / top-$5$ movie-identity
results are finalized. Headline framing: position-in-movie can be
decoded as a continuous target (Table~\ref{tab:stim-decoding});
binning it into $20$ classes tests whether the representation
encodes coarser temporal context as a categorical signal.
Reference Table~\ref{tab:movie-id}.}

% =====================================================================
\subsection{Ablations}
\label{sec:results:ablations}
\todo{Fill in. Recommended ablations to present here: (i) Raw
$C{=}129$ vs.\ CorrCA-5 input under the same Cine-JEPA pipeline;
(ii) per-recording vs.\ global Z-score; (iii) VICReg coefficients
$\lambda_\sigma = \lambda_c \in \{0, 0.25, 1.0\}$; (iv) predictor
bottleneck $D' \in \{16, 24, 32, 64\}$; (v) visual-processing
delay $\Delta$ sweep, validating the $300$\,ms choice from
\S\ref{sec:method:data}.}
 from the main file.
% Citation keys are placeholders; resolve them in the bibliography.
% =====================================================================

\section{Results}
\label{sec:results}

We report decoding performance under the protocol of
\S\ref{sec:method:eval}: every method exposes a deterministic
clip-to-feature map and is fed into the same closed-form linear
probe on the same R6 test split. Each headline number is built in
two stages. \emph{Within} each pretraining seed we compute a
recording-level bootstrap mean of the test metric: the $108$ R6
recordings are resampled with replacement $B = 2{,}000$ times,
the metric is recomputed on each resample, and we take the
bootstrap mean as the per-seed score. \emph{Across} pretraining
seeds we then report the mean and one standard deviation of the
per-seed bootstrap means over five encoder seeds (CBraMod: four,
because of compute constraints). A row is marked
``$\checkmark$'' when the per-seed bootstrap means beat their
chance value at $p < 0.05$ under a one-sample two-sided $t$-test
($r = 0$ for Pearson $r$; $0.5$ for AUC and balanced accuracy;
$0.05/0.25$ for top-$1$/top-$5$ over $20$ bins).

% =====================================================================
\subsection{Stimulus-locked decoding}
\label{sec:results:stim}

Table~\ref{tab:stim-decoding} reports Pearson $r$ between probe
predictions and four continuous movie-feature targets — frame
luminance, frame contrast, normalized temporal position in the
movie, and the narrative event score — on the held-out R6 test
split.

% Stimulus-locked decoding table — referenced from results.tex
% Requires: booktabs
\begin{table}[t]
\centering
\small
\caption{Stimulus-locked decoding on the held-out R6 test split.
All values are Pearson $r$, mean $\pm 1\sigma$ across $5$ encoder
seeds (CBraMod: $n{=}4$), after a recording-level bootstrap
($B{=}2{,}000$) within each seed. ``$\checkmark$'' marks rows
significant against $r{=}0$ at $p{<}0.05$ (one-sample two-sided
$t$-test on the per-seed bootstrap means); ``$n.s.$'' marks
non-significance. Best per column in \textbf{bold}.}
\label{tab:stim-decoding}
\begin{tabular}{lcccc}
\toprule
Method & luminance & contrast & position & narrative \\
\midrule
chance & $0.000$ & $0.000$ & $0.000$ & $0.000$ \\
\addlinespace[2pt]
\multicolumn{5}{l}{\textit{Trivial baselines (no encoder)}} \\
Raw $129$-ch stats & $-0.016 \pm 0.005\,n.s.$ & $-0.008 \pm 0.022\,n.s.$ & $+0.006 \pm 0.025\,n.s.$ & $-0.006 \pm 0.015\,n.s.$ \\
CorrCA stats       & $0.167 \pm 0.041\,\checkmark$ & $0.177 \pm 0.043\,\checkmark$ & $0.158 \pm 0.032\,\checkmark$ & $0.040 \pm 0.023\,\checkmark$ \\
\addlinespace[2pt]
\multicolumn{5}{l}{\textit{Linear / supervised baselines}} \\
mTRF ($129$ channels) & $0.052 \pm 0.033\,\checkmark$ & $0.044 \pm 0.040\,n.s.$ & $0.027 \pm 0.037\,n.s.$ & $0.032 \pm 0.014\,\checkmark$ \\
mTRF (CorrCA-$5$)     & $0.117 \pm 0.011\,\checkmark$ & $0.052 \pm 0.026\,\checkmark$ & $0.071 \pm 0.032\,\checkmark$ & $0.038 \pm 0.040\,n.s.$ \\
Supervised end-to-end (Deep4Net) & $0.207 \pm 0.013\,\checkmark$ & $0.155 \pm 0.017\,\checkmark$ & $0.146 \pm 0.030\,\checkmark$ & $0.144 \pm 0.025\,\checkmark$ \\
\addlinespace[2pt]
\multicolumn{5}{l}{\textit{EEG foundation models, linear probe (frozen)}} \\
BENDR  & --- & --- & --- & --- \\
BIOT   & $0.127 \pm 0.033\,\checkmark$ & $0.107 \pm 0.035\,\checkmark$ & $0.107 \pm 0.038\,\checkmark$ & $0.003 \pm 0.009\,n.s.$ \\
LaBraM & --- & --- & --- & --- \\
Luna   & $0.182 \pm 0.021\,\checkmark$ & $0.158 \pm 0.018\,\checkmark$ & $0.171 \pm 0.012\,\checkmark$ & $0.053 \pm 0.036\,\checkmark$ \\
\addlinespace[2pt]
\multicolumn{5}{l}{\textit{EEG foundation models, full fine-tune}} \\
BENDR  & --- & --- & --- & --- \\
BIOT   & $0.085 \pm 0.060\,\checkmark$ & $0.074 \pm 0.066\,n.s.$ & $0.109 \pm 0.030\,\checkmark$ & $0.013 \pm 0.028\,n.s.$ \\
LaBraM & --- & --- & --- & --- \\
Luna   & $0.192 \pm 0.120\,\checkmark$ & $0.169 \pm 0.038\,\checkmark$ & $0.184 \pm 0.031\,\checkmark$ & $0.095 \pm 0.100\,n.s.$ \\
Luna + CorrCA-$5$ & $0.159 \pm 0.035\,\checkmark$ & $0.134 \pm 0.059\,\checkmark$ & $0.133 \pm 0.062\,\checkmark$ & $0.058 \pm 0.040\,\checkmark$ \\
CBraMod ($n{=}4$) & $0.052 \pm 0.077\,n.s.$ & $0.034 \pm 0.042\,n.s.$ & $0.022 \pm 0.079\,n.s.$ & $0.040 \pm 0.067\,n.s.$ \\
\addlinespace[2pt]
\textbf{Cine-JEPA (ours)} & $\mathbf{0.226 \pm 0.018}\,\checkmark$ & $\mathbf{0.223 \pm 0.014}\,\checkmark$ & $\mathbf{0.226 \pm 0.016}\,\checkmark$ & $\mathbf{0.155 \pm 0.023}\,\checkmark$ \\
\bottomrule
\end{tabular}
\end{table}


\paragraph{Cine-JEPA tops every target.}
Cine-JEPA reaches Pearson $r = 0.226$ on luminance, $0.223$ on
contrast, $0.226$ on temporal position, and $0.155$ on the
narrative event score — the highest score in every column. Its
margin over the strongest non-Cine-JEPA method ranges from $9\%$
relative on luminance (vs.\ supervised Deep4Net at $0.207$) to
$26\%$ on contrast (vs.\ CorrCA-stats at $0.177$) and $23\%$ on
temporal position (vs.\ Luna full-finetune at $0.184$); on
narrative it is the only method to clear $r = 0.1$.

\paragraph{Without a stimulus-locked spatial filter, linear features carry no signal.}
Seven simple summary statistics (mean, std, log-power in
$\delta\!/\!\theta\!/\!\alpha\!/\!\beta\!/\!\gamma$ bands)
computed on the raw $129$-electrode input are statistically
indistinguishable from chance on every target ($|r| \le 0.016$,
all $n.s.$). The same statistics computed on the five
CorrCA-projected channels jump to $r \approx 0.16$--$0.18$ on
luminance, contrast, and temporal position. CorrCA is therefore
doing most of the heavy lifting that any \emph{linear} method can
extract from the EEG. It is, however, far from sufficient:
Cine-JEPA improves on CorrCA-stats by $0.04$--$0.07$ on the three
low/mid-level targets and by an absolute $+0.115$ on the
narrative event score, so the encoder is contributing real
nonlinear structure on top of the spatial filter.

\paragraph{The standard naturalistic-EEG decoder lags every learned method.}
Linear TRFs — the field-standard tool for naturalistic EEG
decoding \citep{crosse2016mtrf} — top out at $r = 0.117$ on
luminance under the CorrCA-5 input and decay to chance on
narrative ($r = 0.038$, $n.s.$). The fully-supervised end-to-end
Deep4Net is much stronger ($r = 0.207, 0.155, 0.146, 0.144$
across the four targets) but still trails Cine-JEPA on every
target by $0.011$ to $0.068$.

\paragraph{Foundation-model probes plateau below Cine-JEPA.}
Among the EEG foundation models we evaluated, Luna is the
strongest: a frozen linear probe on Luna features reaches
$r = 0.182, 0.158, 0.171, 0.053$, second only to Cine-JEPA on
three of four targets. Full fine-tuning of Luna gives a small
additional gain on contrast and temporal position but does not
close the gap, and feeding Luna a CorrCA-5 input
\emph{degrades} its scores rather than improves them — Luna's
own input pipeline already extracts a richer projection than
five linear components, so collapsing the channel dimension up
front discards information. BIOT is intermediate, CBraMod is
within noise of chance under our protocol, and BENDR / LaBraM
did not produce usable scores under the protocol and are listed
as ``---''.

\paragraph{Narrative event decoding is the discriminating target.}
All methods take their largest hit on the narrative event score:
Cine-JEPA is the only method to clear $r = 0.1$, and the
runner-up (Deep4Net at $0.144$) is itself supervised end-to-end
on the same target. Because the narrative score is higher-level
and less stimulus-locked than per-frame luminance or contrast,
the gap on this column separates methods that have learned a
temporally structured representation from methods that have
mainly captured low-level visual statistics. We return to this
in the discussion.

% =====================================================================
\subsection{Subject-trait decoding}
\label{sec:results:traits}
\todo{Fill in once the age-regression and sex-classification
results are finalized. Headline framing: do the same frozen
representations that win at stimulus-locked decoding also carry
recoverable subject-level structure (age as continuous Pearson
$r$; sex as ROC-AUC and balanced accuracy)? Reference
Table~\ref{tab:trait-decoding}.}

% =====================================================================
\subsection{Movie identity}
\label{sec:results:movie-id}
\todo{Fill in once the $20$-bin top-$1$ / top-$5$ movie-identity
results are finalized. Headline framing: position-in-movie can be
decoded as a continuous target (Table~\ref{tab:stim-decoding});
binning it into $20$ classes tests whether the representation
encodes coarser temporal context as a categorical signal.
Reference Table~\ref{tab:movie-id}.}

% =====================================================================
\subsection{Ablations}
\label{sec:results:ablations}
\todo{Fill in. Recommended ablations to present here: (i) Raw
$C{=}129$ vs.\ CorrCA-5 input under the same Cine-JEPA pipeline;
(ii) per-recording vs.\ global Z-score; (iii) VICReg coefficients
$\lambda_\sigma = \lambda_c \in \{0, 0.25, 1.0\}$; (iv) predictor
bottleneck $D' \in \{16, 24, 32, 64\}$; (v) visual-processing
delay $\Delta$ sweep, validating the $300$\,ms choice from
\S\ref{sec:method:data}.}
 from the main file.
% Citation keys are placeholders; resolve them in the bibliography.
% =====================================================================

\section{Results}
\label{sec:results}

We report decoding performance under the protocol of
\S\ref{sec:method:eval}: every method exposes a deterministic
clip-to-feature map and is fed into the same closed-form linear
probe on the same R6 test split. Each headline number is built in
two stages. \emph{Within} each pretraining seed we compute a
recording-level bootstrap mean of the test metric: the $108$ R6
recordings are resampled with replacement $B = 2{,}000$ times,
the metric is recomputed on each resample, and we take the
bootstrap mean as the per-seed score. \emph{Across} pretraining
seeds we then report the mean and one standard deviation of the
per-seed bootstrap means over five encoder seeds (CBraMod: four,
because of compute constraints). A row is marked
``$\checkmark$'' when the per-seed bootstrap means beat their
chance value at $p < 0.05$ under a one-sample two-sided $t$-test
($r = 0$ for Pearson $r$; $0.5$ for AUC and balanced accuracy;
$0.05/0.25$ for top-$1$/top-$5$ over $20$ bins).

% =====================================================================
\subsection{Stimulus-locked decoding}
\label{sec:results:stim}

Table~\ref{tab:stim-decoding} reports Pearson $r$ between probe
predictions and four continuous movie-feature targets — frame
luminance, frame contrast, normalized temporal position in the
movie, and the narrative event score — on the held-out R6 test
split.

% Stimulus-locked decoding table — referenced from results.tex
% Requires: booktabs
\begin{table}[t]
\centering
\small
\caption{Stimulus-locked decoding on the held-out R6 test split.
All values are Pearson $r$, mean $\pm 1\sigma$ across $5$ encoder
seeds (CBraMod: $n{=}4$), after a recording-level bootstrap
($B{=}2{,}000$) within each seed. ``$\checkmark$'' marks rows
significant against $r{=}0$ at $p{<}0.05$ (one-sample two-sided
$t$-test on the per-seed bootstrap means); ``$n.s.$'' marks
non-significance. Best per column in \textbf{bold}.}
\label{tab:stim-decoding}
\begin{tabular}{lcccc}
\toprule
Method & luminance & contrast & position & narrative \\
\midrule
chance & $0.000$ & $0.000$ & $0.000$ & $0.000$ \\
\addlinespace[2pt]
\multicolumn{5}{l}{\textit{Trivial baselines (no encoder)}} \\
Raw $129$-ch stats & $-0.016 \pm 0.005\,n.s.$ & $-0.008 \pm 0.022\,n.s.$ & $+0.006 \pm 0.025\,n.s.$ & $-0.006 \pm 0.015\,n.s.$ \\
CorrCA stats       & $0.167 \pm 0.041\,\checkmark$ & $0.177 \pm 0.043\,\checkmark$ & $0.158 \pm 0.032\,\checkmark$ & $0.040 \pm 0.023\,\checkmark$ \\
\addlinespace[2pt]
\multicolumn{5}{l}{\textit{Linear / supervised baselines}} \\
mTRF ($129$ channels) & $0.052 \pm 0.033\,\checkmark$ & $0.044 \pm 0.040\,n.s.$ & $0.027 \pm 0.037\,n.s.$ & $0.032 \pm 0.014\,\checkmark$ \\
mTRF (CorrCA-$5$)     & $0.117 \pm 0.011\,\checkmark$ & $0.052 \pm 0.026\,\checkmark$ & $0.071 \pm 0.032\,\checkmark$ & $0.038 \pm 0.040\,n.s.$ \\
Supervised end-to-end (Deep4Net) & $0.207 \pm 0.013\,\checkmark$ & $0.155 \pm 0.017\,\checkmark$ & $0.146 \pm 0.030\,\checkmark$ & $0.144 \pm 0.025\,\checkmark$ \\
\addlinespace[2pt]
\multicolumn{5}{l}{\textit{EEG foundation models, linear probe (frozen)}} \\
BENDR  & --- & --- & --- & --- \\
BIOT   & $0.127 \pm 0.033\,\checkmark$ & $0.107 \pm 0.035\,\checkmark$ & $0.107 \pm 0.038\,\checkmark$ & $0.003 \pm 0.009\,n.s.$ \\
LaBraM & --- & --- & --- & --- \\
Luna   & $0.182 \pm 0.021\,\checkmark$ & $0.158 \pm 0.018\,\checkmark$ & $0.171 \pm 0.012\,\checkmark$ & $0.053 \pm 0.036\,\checkmark$ \\
\addlinespace[2pt]
\multicolumn{5}{l}{\textit{EEG foundation models, full fine-tune}} \\
BENDR  & --- & --- & --- & --- \\
BIOT   & $0.085 \pm 0.060\,\checkmark$ & $0.074 \pm 0.066\,n.s.$ & $0.109 \pm 0.030\,\checkmark$ & $0.013 \pm 0.028\,n.s.$ \\
LaBraM & --- & --- & --- & --- \\
Luna   & $0.192 \pm 0.120\,\checkmark$ & $0.169 \pm 0.038\,\checkmark$ & $0.184 \pm 0.031\,\checkmark$ & $0.095 \pm 0.100\,n.s.$ \\
Luna + CorrCA-$5$ & $0.159 \pm 0.035\,\checkmark$ & $0.134 \pm 0.059\,\checkmark$ & $0.133 \pm 0.062\,\checkmark$ & $0.058 \pm 0.040\,\checkmark$ \\
CBraMod ($n{=}4$) & $0.052 \pm 0.077\,n.s.$ & $0.034 \pm 0.042\,n.s.$ & $0.022 \pm 0.079\,n.s.$ & $0.040 \pm 0.067\,n.s.$ \\
\addlinespace[2pt]
\textbf{Cine-JEPA (ours)} & $\mathbf{0.226 \pm 0.018}\,\checkmark$ & $\mathbf{0.223 \pm 0.014}\,\checkmark$ & $\mathbf{0.226 \pm 0.016}\,\checkmark$ & $\mathbf{0.155 \pm 0.023}\,\checkmark$ \\
\bottomrule
\end{tabular}
\end{table}


\paragraph{Cine-JEPA tops every target.}
Cine-JEPA reaches Pearson $r = 0.226$ on luminance, $0.223$ on
contrast, $0.226$ on temporal position, and $0.155$ on the
narrative event score — the highest score in every column. Its
margin over the strongest non-Cine-JEPA method ranges from $9\%$
relative on luminance (vs.\ supervised Deep4Net at $0.207$) to
$26\%$ on contrast (vs.\ CorrCA-stats at $0.177$) and $23\%$ on
temporal position (vs.\ Luna full-finetune at $0.184$); on
narrative it is the only method to clear $r = 0.1$.

\paragraph{Without a stimulus-locked spatial filter, linear features carry no signal.}
Seven simple summary statistics (mean, std, log-power in
$\delta\!/\!\theta\!/\!\alpha\!/\!\beta\!/\!\gamma$ bands)
computed on the raw $129$-electrode input are statistically
indistinguishable from chance on every target ($|r| \le 0.016$,
all $n.s.$). The same statistics computed on the five
CorrCA-projected channels jump to $r \approx 0.16$--$0.18$ on
luminance, contrast, and temporal position. CorrCA is therefore
doing most of the heavy lifting that any \emph{linear} method can
extract from the EEG. It is, however, far from sufficient:
Cine-JEPA improves on CorrCA-stats by $0.04$--$0.07$ on the three
low/mid-level targets and by an absolute $+0.115$ on the
narrative event score, so the encoder is contributing real
nonlinear structure on top of the spatial filter.

\paragraph{The standard naturalistic-EEG decoder lags every learned method.}
Linear TRFs — the field-standard tool for naturalistic EEG
decoding \citep{crosse2016mtrf} — top out at $r = 0.117$ on
luminance under the CorrCA-5 input and decay to chance on
narrative ($r = 0.038$, $n.s.$). The fully-supervised end-to-end
Deep4Net is much stronger ($r = 0.207, 0.155, 0.146, 0.144$
across the four targets) but still trails Cine-JEPA on every
target by $0.011$ to $0.068$.

\paragraph{Foundation-model probes plateau below Cine-JEPA.}
Among the EEG foundation models we evaluated, Luna is the
strongest: a frozen linear probe on Luna features reaches
$r = 0.182, 0.158, 0.171, 0.053$, second only to Cine-JEPA on
three of four targets. Full fine-tuning of Luna gives a small
additional gain on contrast and temporal position but does not
close the gap, and feeding Luna a CorrCA-5 input
\emph{degrades} its scores rather than improves them — Luna's
own input pipeline already extracts a richer projection than
five linear components, so collapsing the channel dimension up
front discards information. BIOT is intermediate, CBraMod is
within noise of chance under our protocol, and BENDR / LaBraM
did not produce usable scores under the protocol and are listed
as ``---''.

\paragraph{Narrative event decoding is the discriminating target.}
All methods take their largest hit on the narrative event score:
Cine-JEPA is the only method to clear $r = 0.1$, and the
runner-up (Deep4Net at $0.144$) is itself supervised end-to-end
on the same target. Because the narrative score is higher-level
and less stimulus-locked than per-frame luminance or contrast,
the gap on this column separates methods that have learned a
temporally structured representation from methods that have
mainly captured low-level visual statistics. We return to this
in the discussion.

% =====================================================================
\subsection{Subject-trait decoding}
\label{sec:results:traits}
\todo{Fill in once the age-regression and sex-classification
results are finalized. Headline framing: do the same frozen
representations that win at stimulus-locked decoding also carry
recoverable subject-level structure (age as continuous Pearson
$r$; sex as ROC-AUC and balanced accuracy)? Reference
Table~\ref{tab:trait-decoding}.}

% =====================================================================
\subsection{Movie identity}
\label{sec:results:movie-id}
\todo{Fill in once the $20$-bin top-$1$ / top-$5$ movie-identity
results are finalized. Headline framing: position-in-movie can be
decoded as a continuous target (Table~\ref{tab:stim-decoding});
binning it into $20$ classes tests whether the representation
encodes coarser temporal context as a categorical signal.
Reference Table~\ref{tab:movie-id}.}

% =====================================================================
\subsection{Ablations}
\label{sec:results:ablations}
\todo{Fill in. Recommended ablations to present here: (i) Raw
$C{=}129$ vs.\ CorrCA-5 input under the same Cine-JEPA pipeline;
(ii) per-recording vs.\ global Z-score; (iii) VICReg coefficients
$\lambda_\sigma = \lambda_c \in \{0, 0.25, 1.0\}$; (iv) predictor
bottleneck $D' \in \{16, 24, 32, 64\}$; (v) visual-processing
delay $\Delta$ sweep, validating the $300$\,ms choice from
\S\ref{sec:method:data}.}
